% Book 4, Year 2, chapter 9 -- Functional programming.
% Provenance (never printed): sources/book4/BRIEFS.md, chapter 9.
\chapter{Functional Programming}\label{ch:b2:functional-programming}

The sum of the squares of the odd readings, written twice: as a loop that changes a variable, and as one expression
that says what is computed.

\omcode{code/bachelor-2/09-functional-programming/listings/odd_squares.rs}{3}{14}{A loop and a
pipeline}\label{lst:b2:functional-programming:odd}

\omoutput{code/bachelor-2/09-functional-programming/outputs/odd_squares.txt}{$7^2 + 3^2 + 9^2 + 15^2 = 364$, twice.}

The second form reads like a spreadsheet formula, \emph{sum the squares of the odd cells}: no counter, no variable
that changes, no order of steps to follow. Functional programming builds programs from such expressions, functions
passed to functions, and data that is never modified. This chapter shows the ideas in Rust, which \omterm{def:b2:rust-ownership:borrow}{borrows} them
freely, and what each costs.

\begin{recall}
Book 2: a recursive function calls itself on a smaller input, and each pending call occupies a frame of the stack.
Book 3: lists and trees are inductive sets, and a property of all of them is proved by structural induction.
\Cref{ch:b2:rust-ownership,ch:b2:rust-enums-and-errors,ch:b2:rust-traits-and-generics}: \omterm{def:b2:rust-ownership:ownership}{ownership} and borrowing,
recursive enums with \code{Box}, \omterm{def:b2:rust-traits-and-generics:trait}{traits} and \omterm{def:b2:rust-traits-and-generics:trait-object}{trait objects}.
\end{recall}

\section{Functions as values}\label{sec:b2:functional-programming:closures}

\begin{definition}[First-class function, higher-order function]\label{def:b2:functional-programming:higher-order}
Functions are \emph{first-class}\index{first-class function} in a language if they are values: stored in variables
and data structures, passed as arguments, returned as results. A \emph{higher-order function}\index{higher-order
function} takes a function as argument or returns one.
\end{definition}

\begin{definition}[Closure, captured environment]\label{def:b2:functional-programming:closure}
A \emph{closure}\index{closure} is a function value written inline, \code{|x| x * rate}, that may use variables of
the scope where it is written. Those variables form its \emph{captured environment}\index{captured environment}: the
closure stores a reference to each, a mutable reference, or the value itself (\code{move}). By what its body does with
them, a Rust closure implements \code{Fn} (it only reads them, and may be called any number of times), \code{FnMut}
(it modifies them: called through a mutable \omterm{def:b2:rust-ownership:borrow}{borrow}), or only \code{FnOnce} (it gives a captured value away: callable
once).
\end{definition}

\omcode{code/bachelor-2/09-functional-programming/listings/closures.rs}{2}{30}{Three kinds of capture, and \omterm{def:b2:functional-programming:closure}{closures}
returned}\label{lst:b2:functional-programming:closures}

\omoutput{code/bachelor-2/09-functional-programming/outputs/closures.txt}{The fines of 30 days with 21 and 14 days of
grace.}\label{out:b2:functional-programming:closures}

\code{make_fine} returns a \omterm{def:b2:functional-programming:closure}{closure}: its parameters die with the call, so the \omterm{def:b2:functional-programming:closure}{closure} must own copies, hence
\code{move}. Each \omterm{def:b2:functional-programming:closure}{closure} has its own anonymous type, which is why \code{make_fine} declares its result as
\code{impl Fn(u32) -> u32}, and why \omterm{def:b2:functional-programming:closure}{closures} of different bodies can share a vector only as \code{Box<dyn Fn>},
\omterm{def:b2:rust-traits-and-generics:trait-object}{trait objects} (\cref{ch:b2:rust-traits-and-generics}). A \omterm{def:b2:functional-programming:closure}{closure} that gives its capture away cannot run twice:

\omoutput[2-7,9-13]{code/bachelor-2/09-functional-programming/outputs/fn_once_error.txt}{\code{fn_once_error.rs} calls
\code{give} twice.}\label{out:b2:functional-programming:once}

\begin{omfigure}
\begin{tikzpicture}[font=\scriptsize, c/.style={draw, minimum width=19mm, minimum height=5mm}]
\node[font=\scriptsize\bfseries] at (1,1.0) {\texttt{fine}: borrows};
\node[c, fill=omDef!10] (f1) at (1,0.2) {\texttt{\&rate}};
\node[text=black!60] at (1,-0.45) {8 bytes, \texttt{Fn}};
\node[c, fill=omThm!10] (r) at (1,-1.6) {\texttt{rate = 10}};
\node[text=black!60, anchor=north] at (1,-1.9) {in \texttt{main}'s frame};
\draw[ptr] (f1.west) to[out=200, in=160] (r.west);
\begin{scope}[xshift=5cm]
\node[font=\scriptsize\bfseries] at (1,1.0) {\texttt{count}: borrows mutably};
\node[c, fill=omDef!10] (f2) at (1,0.2) {\texttt{\&mut calls}};
\node[text=black!60] at (1,-0.45) {8 bytes, \texttt{FnMut}};
\node[c, fill=omThm!10] (c) at (1,-1.6) {\texttt{calls = 0}, then 2};
\draw[ptr] (f2.west) to[out=200, in=160] (c.west);
\end{scope}
\begin{scope}[xshift=10cm]
\node[font=\scriptsize\bfseries] at (1,1.0) {\texttt{make\_fine(10, 21)}: owns};
\node[c, fill=omDef!10, minimum width=30mm] (f3) at (1,0.2) {\texttt{cents\_per\_day = 10}};
\node[c, fill=omDef!10, minimum width=30mm, below=0pt of f3] {\texttt{grace = 21}};
\node[text=black!60] at (1,-0.75) {8 bytes, \texttt{Fn}, outlives the call};
\end{scope}
\end{tikzpicture}
\par\smallskip{\small A \omterm{def:b2:functional-programming:closure}{closure} is a struct of its captures plus a function that receives that struct. The first two
hold references into \code{main}'s frame and may not outlive it; the third holds its own copies.}
\end{omfigure}

\section{Map, filter, fold and lazy iterators}\label{sec:b2:functional-programming:fold}

\begin{definition}[Map operation, filter, fold]\label{def:b2:functional-programming:fold}
On a sequence $[x_1, \dots, x_n]$, the \emph{map operation}\index{map operation} with a function $f$ gives $[f(x_1),
\dots, f(x_n)]$; the \emph{filter}\index{filter} with a predicate $p$ keeps the $x_i$ with $p(x_i)$ true, in order; the
\emph{fold}\index{fold} from $z$ with a two-argument function $g$ gives $g(\dots g(g(z, x_1), x_2) \dots, x_n)$ (a left
fold; the right fold is $g(x_1, g(x_2, \dots g(x_n, z)))$). Sums, counts, maxima, and the building of collections are
folds.
\end{definition}

\begin{definition}[Lazy iterator, iterator adaptor]\label{def:b2:functional-programming:lazy}
A Rust iterator is \emph{lazy}\index{lazy iterator}: it computes the next element only when \code{next} is called.
An \emph{iterator adaptor}\index{iterator adaptor} (\code{map}, \code{filter}, \code{take}, \dots) wraps an iterator
into a new one and does no work when created; a consumer (\code{sum}, \code{collect}, \code{count}, a \code{for} loop)
calls \code{next} until the end or until it has enough.
\end{definition}

\omcode{code/bachelor-2/09-functional-programming/listings/lazy.rs}{3}{16}{\omterm{def:b2:functional-programming:pure}{Side effects} that show when each step
runs}\label{lst:b2:functional-programming:lazy}

\omoutput{code/bachelor-2/09-functional-programming/outputs/lazy.txt}{One element at a time, and the last page
count is never read.}\label{out:b2:functional-programming:lazy}

\begin{omfigure}
\begin{tikzpicture}[font=\scriptsize, s/.style={draw, rounded corners=2pt, fill=omDef!10, minimum width=20mm, minimum height=7mm, align=center}]
\node[s, fill=black!5] (src) at (0,0) {\texttt{pages.iter()}};
\node[s] (m) at (3.2,0) {\texttt{map}\\ \texttt{p * 2}};
\node[s] (f) at (6.4,0) {\texttt{filter}\\ \texttt{d > 300}};
\node[s] (t) at (9.6,0) {\texttt{take(2)}};
\node[s, fill=omThm!10] (c) at (12.6,0) {\texttt{collect}};
\foreach \a/\b in {c/t, t/f, f/m, m/src}
  \draw[ptr, dashed] ([yshift=2mm]\a.west) -- node[above, text=black!60] {\texttt{next}} ([yshift=2mm]\b.east);
\foreach \a/\b in {src/m, m/f, f/t, t/c}
  \draw[ptr] ([yshift=-2mm]\a.east) -- ([yshift=-2mm]\b.west);
\node[text=black!60, align=center] at (4.8,-1.0) {65 $\to$ 130, refused; 133 $\to$ 266, refused;\\ 412 $\to$ 824, passed up};
\end{tikzpicture}
\par\smallskip{\small A \omterm{def:b2:functional-programming:lazy}{lazy} pipeline is pulled by its consumer. \code{collect} asks \code{take} for an element,
which asks \code{filter}, which asks \code{map} until one passes; after two, \code{take} stops asking, and 320 is
never mapped.}
\end{omfigure}

Laziness makes infinite sequences usable, as long as the consumer stops:

\omcode{code/bachelor-2/09-functional-programming/listings/infinite.rs}{3}{17}{Three infinite
iterators}\label{lst:b2:functional-programming:infinite}

\omoutput{code/bachelor-2/09-functional-programming/outputs/infinite.txt}{Each sequence is computed only as far as
asked.}\label{out:b2:functional-programming:infinite}

On real data, a pipeline states the cleaning rules one per line. The weather log of
\cref{ch:b2:rust-enums-and-errors} gives each station's mean temperature:

\omcode{code/bachelor-2/09-functional-programming/listings/pipeline.rs}{5}{12}{A mean as a
\omterm{def:b2:functional-programming:fold}{fold}}\label{lst:b2:functional-programming:pipeline}

\omoutput[2-4]{code/bachelor-2/09-functional-programming/outputs/run_pipeline.txt}{The duplicated Riverside line is
counted here, unlike in \cref{ch:b2:rust-enums-and-errors}.}\label{out:b2:functional-programming:pipeline}

The whole pipeline is one loop over the lines, with no intermediate vector: each line is split, tested, parsed and
added before the next is read. Rust compiles such chains to the loop one would write by hand.

\section{Immutable data}\label{sec:b2:functional-programming:persistent}

A functional program does not modify data: an update builds a new value. Done naively, every update copies
everything. Done well, the new value shares all the parts it did not change with the old one.

\begin{definition}[Persistent data structure, structural sharing, path copying]\label{def:b2:functional-programming:persistent}
A data structure is \emph{persistent}\index{persistent data structure} if an update returns a new version and leaves
every old version intact and usable. \emph{Structural sharing}\index{structural sharing} is the use of the same
nodes by several versions; \emph{path copying}\index{path copying} implements an update of a tree by copying only the
nodes on the path from the root to the place of the change, the new nodes pointing to the old, unchanged subtrees.
\end{definition}

Sharing needs \omterm{def:b2:cpp-copy-move:shared-ptr}{shared ownership}: a node with several \omterm{def:b2:rust-ownership:ownership}{owners} must be freed only when the last one goes.
\code{Rc<T>} (reference counted) does this: \code{Rc::clone} copies the pointer and increments a count stored next
to the value; dropping an \code{Rc} decrements it, and the value is freed at zero.

\omcode{code/bachelor-2/09-functional-programming/listings/persistent_list.rs}{5}{36}{Two lists with one
tail}\label{lst:b2:functional-programming:plist}

\omoutput{code/bachelor-2/09-functional-programming/outputs/persistent_list.txt}{The tail is shared, never
copied.}\label{out:b2:functional-programming:plist}

\begin{omfigure}
\begin{tikzpicture}[font=\scriptsize, n/.style={draw, minimum width=8mm, minimum height=5mm, fill=omDef!10}]
\node[n, fill=omThm!10] (a) at (0,0.7) {10};
\node[n, fill=omThm!10] (b) at (0,-0.7) {20};
\node[n] (t3) at (2.2,0) {3};
\node[n] (t2) at (3.8,0) {2};
\node[n] (t1) at (5.4,0) {1};
\node[n, fill=black!5] (nil) at (7.0,0) {Nil};
\draw[ptr] (a.east) -- (t3.west); \draw[ptr] (b.east) -- (t3.west);
\draw[ptr] (t3) -- (t2); \draw[ptr] (t2) -- (t1); \draw[ptr] (t1) -- (nil);
\node[anchor=east] at (-0.6,0.7) {\texttt{a}}; \node[anchor=east] at (-0.6,-0.7) {\texttt{b}};
\node[above=1mm of t3, text=black!60] {count 3: \texttt{tail}, \texttt{a}, \texttt{b}};
\end{tikzpicture}
\par\smallskip{\small Two \omterm{def:b2:functional-programming:persistent}{persistent} lists sharing a tail. Adding an element at the front costs one node, whatever
the length; neither list can change the shared part, since nothing modifies a node once built.}
\end{omfigure}

A search tree is shared the same way. Inserting into a \omterm{def:b2:functional-programming:persistent}{persistent} tree rebuilds the path to the new leaf and points
everything else at the old version:

\omcode{code/bachelor-2/09-functional-programming/listings/persistent_bst.rs}{13}{26}{Insertion by \omterm{def:b2:functional-programming:persistent}{path
copying}}\label{lst:b2:functional-programming:pbst}

\omoutput{code/bachelor-2/09-functional-programming/outputs/persistent_bst.txt}{Inserting 140: four new nodes, four
shared; \code{v1} still has its seven keys.}\label{out:b2:functional-programming:pbst}

\begin{omfigure}
\begin{tikzpicture}[font=\scriptsize, o/.style={draw, rounded corners=2pt, fill=omDef!10, minimum width=9mm, minimum height=5mm}, c/.style={o, fill=omThm!15, thick}]
\node[o] (r1) at (0,0) {189};
\node[o] (a1) at (-1.8,-1) {65};
\node[o] (b1) at (1.8,-1) {251};
\node[o] (a2) at (-2.6,-2) {98};
\node[o] (a3) at (-1.0,-2) {133};
\node[o] (b2) at (2.6,-2) {412};
\node[o] (b3) at (2.0,-3) {320};
\draw[ptr] (r1) -- (a1); \draw[ptr] (r1) -- (b1); \draw[ptr] (a1) -- (a2); \draw[ptr] (a1) -- (a3);
\draw[ptr] (b1) -- (b2); \draw[ptr] (b2) -- (b3);
\node[above=1mm of r1] {\texttt{v1}};
\node[c] (r2) at (-5.4,0) {189};
\node[c] (c1) at (-6.2,-1) {65};
\node[c] (c3) at (-5.4,-2) {133};
\node[c] (new) at (-4.8,-3) {140};
\node[above=1mm of r2] {\texttt{v2}};
\draw[ptr, omThm] (r2) -- (c1); \draw[ptr, omThm] (c1) -- (c3); \draw[ptr, omThm] (c3) -- (new);
\draw[ptr, omThm, dashed] (r2.north) to[out=50, in=110] (b1.north);
\draw[ptr, omThm, dashed] (c1.east) -- (a2.west);
\end{tikzpicture}
\par\smallskip{\small \omterm{def:b2:functional-programming:persistent}{Path copying}, from \cref{out:b2:functional-programming:pbst}. \code{v2} owns three copies (189,
65, 133) and the new 140; its dashed pointers lead into \code{v1}'s nodes 251 and 98, with all that hangs below
them. \code{v1} is untouched.}
\end{omfigure}

\begin{proposition}[Path copying]\label{prop:b2:functional-programming:path-copy}
Inserting a key absent from a \omterm{def:b2:functional-programming:persistent}{persistent} search tree, at depth $d$ (the root at depth 0), creates exactly $d + 1$
nodes: copies of the $d$ nodes on the path and the new leaf. Every other node of the old version is shared by the new
one.
\end{proposition}

\begin{proof}
By induction on $d$. If $d = 0$, the tree is empty and \code{insert} creates one node. If $d \ge 1$, the root is
not the key; \code{insert} creates one copy of the root and calls itself on one subtree, where the key will be at
depth $d - 1$: by induction that call creates $d$ nodes and shares every other node of that subtree. The copy of the
root points to the other subtree unchanged (\code{n.right.clone()} or \code{n.left.clone()}, a pointer copy). In
all, $1 + d$ new nodes, and every old node off the path is reachable from the new root.
\end{proof}

\section{Recursion and tail calls}\label{sec:b2:functional-programming:tail}

Without loops that update variables, functional programs repeat by recursion. Each pending call costs a stack frame.

\begin{definition}[Tail call, tail recursion]\label{def:b2:functional-programming:tail}
A call is a \emph{tail call}\index{tail call} if it is the last action of the calling function: its result is
returned as it is, with nothing left to compute. A function is \emph{tail-recursive}\index{tail recursion} if its
recursive calls are tail calls. A compiler that performs \emph{tail-call elimination} turns a tail call into a jump,
reusing the frame.
\end{definition}

\omcode{code/bachelor-2/09-functional-programming/listings/deep_recursion.rs}{3}{15}{A sum with one frame per
element}\label{lst:b2:functional-programming:deep}

\omoutput{code/bachelor-2/09-functional-programming/outputs/run_deep.txt}{A million frames do not fit in 1\,MiB, with
or without optimisation. Status 134 is the abort signal: $128 + 6$.}\label{out:b2:functional-programming:deep}

\code{first + sum_rec(rest)} is not a \omterm{def:b2:functional-programming:tail}{tail call}: the addition waits for the result. Rust does not promise \omterm{def:b2:functional-programming:tail}{tail-call
elimination} even for real \omterm{def:b2:functional-programming:tail}{tail calls}; the optimiser may do it, but a program that depends on it is fragile. The
remedy is a method.

\begin{method}[The accumulator transformation]\label{met:b2:functional-programming:accumulator}
To make a recursion that combines its result with pending work into a loop:
\begin{enumerate}
\item add a parameter, the accumulator, that carries the work done so far, and start it at the neutral value of the
combination (0 for a sum);
\item in the recursive case, do the pending work on the accumulator \emph{before} the call, so that the call becomes a
\omterm{def:b2:functional-programming:tail}{tail call}: \code{sum_acc(rest, acc + first)};
\item return the accumulator in the base case;
\item replace the \omterm{def:b2:functional-programming:tail}{tail call} by updating the parameters and jumping back to the start: a \code{loop} or
\code{while}.
\end{enumerate}
The transformation is correct when the combination is associative, since it regroups $x_1 + (x_2 + (\dots))$ as
$((0 + x_1) + x_2) + \dots$.
\end{method}

\omcode{code/bachelor-2/09-functional-programming/listings/tail_loop.rs}{2}{17}{Steps 1--3, then step
4}\label{lst:b2:functional-programming:tail}

\omoutput{code/bachelor-2/09-functional-programming/outputs/tail_loop.txt}{The loop sums the million numbers on the
same 1\,MiB stack; the accumulator version is run on the first thousand.}\label{out:b2:functional-programming:tail}

\section{Reasoning by equations}\label{sec:b2:functional-programming:equations}

\begin{definition}[Pure function, side effect, referential transparency]\label{def:b2:functional-programming:pure}
A \emph{side effect}\index{side effect} is any effect of a call besides its result: a change to a variable outside
it, an input or output, a \omterm{def:b2:rust-enums-and-errors:panic}{panic}. A function is \emph{pure}\index{pure function} if its result depends only on its
arguments and it has no side effect. An expression is \emph{referentially transparent}\index{referential
transparency} if replacing it by its value does not change the program.
\end{definition}

\begin{definition}[Equational reasoning]\label{def:b2:functional-programming:equational}
\emph{Equational reasoning}\index{equational reasoning} proves properties of \omterm{def:b2:functional-programming:pure}{pure} programs by rewriting: each
defining equation of a function may be used, left to right or right to left, as an equality of values, as in
algebra.
\end{definition}

Write a list as $[\,]$ or $x : xs$, an element in front of a list. The \omterm{def:b2:functional-programming:fold}{map operation} and the right \omterm{def:b2:functional-programming:fold}{fold} are defined by
the equations
\begin{align*}
\mathrm{map}\,f\,[\,] &= [\,], & \mathrm{map}\,f\,(x : xs) &= f\,x : \mathrm{map}\,f\,xs,\\
\mathrm{foldr}\,g\,z\,[\,] &= z, & \mathrm{foldr}\,g\,z\,(x : xs) &= g\,x\,(\mathrm{foldr}\,g\,z\,xs).
\end{align*}

\begin{theorem}[Map fusion]\label{thm:b2:functional-programming:map-fusion}
For all functions $f$, $g$ and every finite list $xs$, $\mathrm{map}\,f\,(\mathrm{map}\,g\,xs) = \mathrm{map}\,(f
\circ g)\,xs$.
\end{theorem}

\begin{proof}
By structural induction on $xs$. For $[\,]$: both sides are $[\,]$ by the first equation. For $x : xs$, assuming the
claim for $xs$:
\begin{align*}
\mathrm{map}\,f\,(\mathrm{map}\,g\,(x : xs)) &= \mathrm{map}\,f\,(g\,x : \mathrm{map}\,g\,xs)
= f\,(g\,x) : \mathrm{map}\,f\,(\mathrm{map}\,g\,xs)\\
&= (f \circ g)\,x : \mathrm{map}\,(f \circ g)\,xs,
\end{align*}
by the second equation twice and the hypothesis; the right-hand side is $\mathrm{map}\,(f \circ g)\,(x : xs)$.
\end{proof}

For Rust this says that \code{.map(g).map(f)} and \code{.map(|x| f(g(x)))} give the same elements, provided
\code{f} and \code{g} are \omterm{def:b2:functional-programming:pure}{pure}: with \omterm{def:b2:functional-programming:pure}{side effects} the order of the effects changes (\cref{out:b2:functional-programming:lazy}
shows them interleaved). Compilers use such laws to merge passes; programmers use them to restructure code
without retesting by hand. The test crate \code{tests/functional} checks this law, \omterm{def:b2:functional-programming:fold}{filter} fusion and the \omterm{def:b2:functional-programming:fold}{fold} laws on
500 generated lists.

\begin{proposition}[Map and filter are folds]\label{prop:b2:functional-programming:fold-cons}
For every finite list $xs$, $\mathrm{map}\,f\,xs = \mathrm{foldr}\,(\lambda x\,acc.\ f\,x : acc)\,[\,]\,xs$, and
$\mathrm{filter}\,p\,xs = \mathrm{foldr}\,(\lambda x\,acc.\ \text{if } p\,x \text{ then } x : acc \text{ else }
acc)\,[\,]\,xs$.
\end{proposition}

\begin{proof}
For the map, by induction on $xs$. For $[\,]$ both sides are $[\,]$. For $x : xs$, the right side is $(\lambda x\,acc.\
f\,x : acc)\,x\,(\mathrm{foldr}\,\dots\,xs) = f\,x : \mathrm{foldr}\,\dots\,xs = f\,x : \mathrm{map}\,f\,xs$ by the
hypothesis, which is $\mathrm{map}\,f\,(x : xs)$. For the \omterm{def:b2:functional-programming:fold}{filter}, with $\mathrm{filter}\,p\,(x : xs) = x :
\mathrm{filter}\,p\,xs$ if $p\,x$ and $\mathrm{filter}\,p\,xs$ otherwise, the same induction with two cases.
\end{proof}

\begin{inpractice}[Iterators in practice]
Prefer a chain of adaptors to an index loop when each step has a name (\code{filter}, \code{map}, \code{sum}): it
cannot read out of bounds, and it is as fast. Keep the \omterm{def:b2:functional-programming:closure}{closures} \omterm{def:b2:functional-programming:pure}{pure}; put \omterm{def:b2:functional-programming:pure}{side effects} at the end of the chain, in a
\code{for} loop or a \code{for_each}. Reach for \code{Rc} and \omterm{def:b2:functional-programming:persistent}{persistent} structures when several versions must live
together (undo, snapshots, concurrent readers); for one version updated in place, a plain \code{Vec} or
\code{HashMap} is simpler and faster.
\end{inpractice}

\section{Exercises}\label{sec:b2:functional-programming:exercises}

\begin{exercise}[$\star$]\label{exo:b2:functional-programming:1}
Predict the lines printed, in order, then check in the solution. How many elements of \code{v} are read?
\omcode{code/bachelor-2/09-functional-programming/listings/exo_lazy_order.rs}{3}{8}{A \omterm{def:b2:functional-programming:lazy}{lazy} chain with
\code{inspect}}
\end{exercise}

\begin{exercise}[$\star$]\label{exo:b2:functional-programming:2}
The program does not compile. Read the error in the solution, fix it in the signature of \code{for_each_loan}, and
explain why the compiler cannot accept \code{Fn} here.
\omcode{code/bachelor-2/09-functional-programming/listings/exo_fnmut.rs}{3}{17}{Counting late loans}
\end{exercise}

\begin{exercise}[$\star$]\label{exo:b2:functional-programming:3}
Write \code{compose(f, g)}, which returns the \omterm{def:b2:functional-programming:closure}{closure} $x \mapsto f(g(x))$, generic over the three types involved.
Then write \code{compose_all}, which composes a vector of boxed functions of one type, the last applied first.
\end{exercise}

\begin{exercise}[$\star$]\label{exo:b2:functional-programming:4}
Predict the three sizes printed, on a target with 8-byte pointers, from what each \omterm{def:b2:functional-programming:closure}{closure} captures and how.
\omcode{code/bachelor-2/09-functional-programming/listings/exo_sizes.rs}{5}{13}{Three \omterm{def:b2:functional-programming:closure}{closures}}
\end{exercise}

\begin{exercise}[$\star\star$]\label{exo:b2:functional-programming:5}
Predict the four values. Which of the four \omterm{def:b2:functional-programming:fold}{folds} would give the same result as a right \omterm{def:b2:functional-programming:fold}{fold} with the same function,
arguments swapped?
\omcode{code/bachelor-2/09-functional-programming/listings/exo_folds.rs}{2}{9}{Four \omterm{def:b2:functional-programming:fold}{folds}}
\end{exercise}

\begin{exercise}[$\star\star$]\label{exo:b2:functional-programming:6}
The program means to print the number of long titles, then the titles. Explain the error, and give two fixes: one
that builds the filtered iterator twice, one that collects it once.
\omcode{code/bachelor-2/09-functional-programming/listings/exo_twice.rs}{3}{10}{An iterator used twice}
\end{exercise}

\begin{exercise}[$\star\star$]\label{exo:b2:functional-programming:7}
Write \code{frequencies(text) -> HashMap<String, u32>} as a single chain ending with a \omterm{def:b2:functional-programming:fold}{fold}, words split at
non-alphanumeric characters and lowercased, and \code{top(text, k)}, the $k$ most frequent words, ties broken
alphabetically.
\end{exercise}

\begin{exercise}[$\star\star$]\label{exo:b2:functional-programming:8}
Prove by structural induction that $\mathrm{foldr}\,(:)\,[\,]\,xs = xs$ for every finite list: the right \omterm{def:b2:functional-programming:fold}{fold} of
the list \omterm{def:b2:cpp-classes-raii:constructor}{constructor} is the identity. What does the left \omterm{def:b2:functional-programming:fold}{fold} of $(\lambda acc\,x.\ x : acc)$ from $[\,]$ give?
\end{exercise}

\begin{exercise}[$\star\star$]\label{exo:b2:functional-programming:9}
Write a \omterm{def:b2:functional-programming:persistent}{persistent} queue from two \omterm{def:b2:functional-programming:persistent}{persistent} lists: \code{push} adds to the back list, \code{pop} takes from the front
list, and when the front is empty, reverses the back into the front. Test first-in first-out order and that an old
version still holds its elements.
\end{exercise}

\begin{exercise}[$\star\star\star$]\label{exo:b2:functional-programming:10}
Prove the \omterm{def:b2:functional-programming:fold}{filter} fusion law, $\mathrm{filter}\,p\,(\mathrm{filter}\,q\,xs) = \mathrm{filter}\,(\lambda x.\ q\,x \wedge
p\,x)\,xs$, by structural induction, then the law $\mathrm{map}\,f\,(\mathrm{filter}\,(p \circ f)\,xs) =
\mathrm{filter}\,p\,(\mathrm{map}\,f\,xs)$. Why do both need $p$, $q$ and $f$ \omterm{def:b2:functional-programming:pure}{pure}?
\end{exercise}

\begin{exercise}[$\star\star\star$]\label{exo:b2:functional-programming:11}
Write a \omterm{def:b2:functional-programming:persistent}{persistent} map from \code{u32} keys to values by \omterm{def:b2:functional-programming:persistent}{path copying}, where inserting an existing key replaces its
value in the new version only. Test it by inserting 1000 random pairs, keeping every version, and checking a sample of
versions against a \code{BTreeMap} snapshot taken at the same step.
\end{exercise}

\begin{exercise}[$\star\star\star$]\label{exo:b2:functional-programming:12}
In the queue of \cref{exo:b2:functional-programming:9}, used one version after the other, each element is \omterm{def:b2:rust-ownership:move}{moved
from} the back to the front once, so $n$ operations cost $O(n)$ in all: $O(1)$ amortised. Show that this bound fails
for \omterm{def:b2:functional-programming:persistent}{persistent} use: give a sequence of operations on a queue of $n$ elements, reusing an old version, in which each of
$m$ pops costs $n$ steps, so $m$ pops cost $mn$. Which assumption of the amortised argument breaks?
\end{exercise}

\section{Problem: A spreadsheet that never forgets}\label{sec:b2:functional-programming:problem}

\begin{problem}[{Weekend problem --- A spreadsheet that never forgets}]\label{pb:b2:functional-programming:1}
A spreadsheet's undo button needs every earlier state of the sheet. Saving a full copy at each edit costs $n$ cells
per edit. This problem builds a sheet whose cells are kept in a \omterm{def:b2:functional-programming:persistent}{persistent} balanced search tree, with formulas as
\omterm{def:b2:functional-programming:closure}{closures} evaluated only when asked, and counts what a thousand versions really cost.

\medskip\noindent\textbf{Part I --- Formulas as closures.}\par\nopagebreak
\omcode{code/bachelor-2/09-functional-programming/listings/sheet.rs}{6}{26}{Cells, nodes and versions}
\begin{enumerate}[beginpenalty=10000]
\item What is the type of a formula? Why does it take a \code{&Sheet} rather than capturing one?
\item Why is a formula an \code{Rc<dyn Fn>} and not a \code{Box<dyn Fn>}? (Look at what \code{set} does to the cells
on its path.)
\item Which \omterm{def:b2:rust-traits-and-generics:trait}{trait} among \code{Fn}, \code{FnMut} and \code{FnOnce} must a formula implement, and why not the others?
\item In \code{sheet_demo.rs}, each column total is \code{move |s| (0..9).map(|r| s.value(10 * r + c)).sum()}. What
does it capture, and why must it be \code{move}?
\item Is a formula a \omterm{def:b2:functional-programming:pure}{pure function} of the sheet? What would break if it modified a captured variable?
\item What does a formula that reads its own cell do when evaluated?
\end{enumerate}

\medskip\noindent\textbf{Part II --- Cells evaluated lazily.}\par\nopagebreak
\omcodes{code/bachelor-2/09-functional-programming/listings/sheet.rs}{50-57,70-80}{Making a version, and reading a
value}
\begin{enumerate}[resume, beginpenalty=10000]
\item How many formulas run when the sheet is built? Check against the output.
\item Reading the total of column 5 runs how many formulas? How many cells are read?
\item The counter \code{evaluations} is a \code{Cell<u32>} behind an \code{Rc}. Why not a plain \code{u32}?
\item After the edits, \code{value(95)} is asked of version 0 and of version 1000. Explain the two answers.
\item A formula's value is recomputed at each request. Describe a cache, and the condition on versions under which a
cached value stays valid.
\item Why is laziness useful when most of a thousand versions are never displayed?
\end{enumerate}

\medskip\noindent\textbf{Part III --- Every version kept.}\par\nopagebreak
\omcode{code/bachelor-2/09-functional-programming/listings/sheet.rs}{38}{48}{Setting a cell by \omterm{def:b2:functional-programming:persistent}{path copying}}
\begin{enumerate}[resume, beginpenalty=10000]
\item \code{set} differs from \code{insert} of \cref{lst:b2:functional-programming:pbst}: the key always exists.
How many nodes does setting a cell at depth $d$ create, by \cref{prop:b2:functional-programming:path-copy}?
\item \code{build} puts the middle key at the root. Prove that a tree built this way from $n$ keys has at most
$\lfloor \log_2 n \rfloor + 1$ levels.
\item What does \code{cell.clone()} cost for a \code{Num}, and for a \code{Formula}?
\item \code{distinct_nodes} walks every version but stops at nodes already seen. Why is its cost proportional to the
number of distinct nodes, plus the number of versions?
\item Which nodes do version 0 and version 1000 share? (Think of the cells never edited, and of their
ancestors.)
\item The edits go to cells $37i \bmod 90$. Why are all 90 number cells edited, and how often each?
\end{enumerate}

\medskip\noindent\textbf{Part IV --- Memory, proved and measured.}\par\nopagebreak
\omoutput[2-8]{code/bachelor-2/09-functional-programming/outputs/run_sheet.txt}{100 cells, 1000 edits, all versions
kept.}
\begin{enumerate}[resume, beginpenalty=10000]
\item Prove that $k$ edits of an $n$-cell sheet built by \code{build} create at most $k(\lfloor \log_2 n \rfloor + 1)$
nodes, so that all versions hold at most $n + k(\lceil \log_2 n \rceil + 1)$ nodes.
\item Compare with the count printed. Why is the count below the bound?
\item Full copies would hold $(k + 1)n$ nodes. What is the ratio here, and how does it grow with $n$?
\item At 40 bytes per node, how much memory do the thousand versions take? Is the \code{Rc} count included in the 40?
\item Old versions that nobody holds are freed. Which line of \code{sheet_demo.rs} keeps them all alive?
\item State the result of the problem.
\end{enumerate}
\end{problem}
